
\input texsis
\paper

\overfullrule=0pt
\def\frac#1#2{#1\over #2}
%
\def\bmatrix#1{\left[ \matrix{#1} \right]}
\def\be{{\beta}}
\def\toone{{t+1}}

%\def\frac#1#2{#1\over #2}
%\magnification=1200
\overfullrule=0pt

\line{Advanced Macro \hfil SMU}
\line{Winter 2019 \hfil Thomas Sargent}

%\centerline{\bf Practice problems}
\centerline{\bf Homework 1}


\medskip

\noindent{\bf Instructions:}  Please feel free to work on this problem in groups of 3 or 4 students. But please hand in your own work. This homework is
due on Friday Jan 18.

\medskip
\noindent  {\bf Muth and Kalman again}

\medskip
\noindent Consider the linear state space system
$$\eqalign{ x_{t+1} & = \rho x_t + c w_{t+1} \cr
            y_t & = x_t + v_t , } \ \leqno(1) $$
where $ \rho \in [0, 1)$ and $c$ are scalars, $\{w_t\}_{t=0}^\infty$ is an i.i.d.\ scalar process with $w_t \sim N(0,1)$,
 $\{v_t\}_{t=0}^\infty $ is an i.i.d. process of measurement errors with $E v_t w_{s+1} = 0$ for all $t\geq 0 $ and $s\geq 0 $ with
 $v_t \sim N(0,R)$, and $x_0$ is a random initial value distributed as $x_0 \sim N(0, {\frac{c^2}{1 - \rho^2}})$. ( Note: Muth (1960)
 considered the limiting case in which $\rho \uparrow  1$.)
 %\auth{Muth, John F.}%
 %\auth{Kalman, Rudolf}%


\medskip

\noindent{\bf a.} Is the univariate process $\{x_t\}_{t=0}^\infty$  Markov?


\medskip
\noindent {\bf b.} Please compute the stationary distribution of the stochastic process $\{x_t\}$ and compare it to the initial distribution
assumed here.

\medskip
\noindent{\bf c.} Is the univariate process $\{y_t\}_{t=0}^\infty$ Markov?

\medskip

\noindent {\bf d.}  Find an {\it innovations representation\/} for $\{y_t\}_{t=-\infty}^\infty$ of the time-invariant form
$$\eqalign{ \hat x_{t+1} & = \rho \hat x_t + k a_t \cr
            y_t & = \hat x_t + a_t,  } \ \leqno(2) $$
where $k$ is a regression coefficient,  $\hat x_t = E x_t | y^{t-1}$, $a_t = y_t - E [y_t | y^{t-1}] $,  $a_t$ is by construction serially uncorrelated and homoskedastic
with $E [a_t| y^{t-1}] = 0$ and $E a_t^2 = \omega$, and $y^{t-1}$ is the semi-infinite history $y_t, y_{t-1}, \ldots$.
Please fill in the blanks in this sentence: ``$k$ is the regression coefficient of $\underline{ \ \ \ \ \ \ \ \ \ \ }$  on $\underline{ \ \ \ \ \ \ \ \ \ \ }$.'' Please give formulas for $k$ and $\omega$.



\medskip

\noindent{\bf e.} Please verify the following formulas:
$$ \eqalign{ \hat x_{t+1} & = k \sum_{j=0}^\infty (\rho - k)^j y_{t-j}  \cr
             y_t & = k \sum_{j=0}^\infty (\rho - k)^j y_{t-j-1} + a_t , }$$
where $a_t $ is orthogonal to $y_{t-j-1}$ for $j \geq 1$ so that the second equation is an {\it autoregressive representation} or {\it autoregression}
for $y_t$. Please verify that $E[ y_{t+1} | y^t] = \hat x_{t+1}$.



\medskip

\noindent{\bf f.} Where $z^t$ denotes a semi-infinite history of a scalar process $z_t$, please verify the following formulas:
$$ \eqalign{ x_t & = c \sum_{j=0}^\infty \rho^j w_{t-j} \cr
             y_t & =   c \sum_{j=0}^\infty \rho^j w_{t-j} + v_t \cr
             E[y_t | w^{t-1}, v^{t-1}] & = c \sum_{j=1}^\infty \rho^j w_{t-j} \cr
             E ( y_t - [y_t | w^{t-1}, v^{t-1}])^2 & = c^2 + R .}$$
Please verify that the one-step ahead prediction error variance    $E ( y_t - [y_t | w^{t-1}, v^{t-1}])^2$ is less than
the one-step ahead prediction error variance $E a_t^2$.  Please explain why.

\medskip
\noindent {\bf g.}  Where $L$ is the lag operator presented in appendix A of chapter 2 of RMT5, %
%appendix \use{appa1} of this chapter,
 please verify the following formula:
$$ a_t = {\frac{1}{1- (\rho-k)L}} \left[ cw_t + (1-\rho L) v_t \right]. $$
Please interpret this formula as showing how the contemporaneous scalar shock $a_t$ reflects ``old news'' in the form of past
$w_t$'s and past $v_t$'s.  Please interpret the discrepancy in variances of one-step ahead prediction errors that you computed
in part {\bf f} in light of this formula.




\bye






\bye
